\section*{Chapitre \ref{ch:g6:rocks-and-minerals} --- Roches et minéraux}

\begin{solution}{exo:g6:rocks-and-minerals:1}
Granite~: ignée. Grès~: sédimentaire. Marbre~: métamorphique. Basalte~:
ignée. Calcaire~: sédimentaire. Ardoise~: métamorphique.
\end{solution}

\begin{solution}{exo:g6:rocks-and-minerals:2}
Le quartz est un \omterm{def:g6:rocks-and-minerals:mineral}{minéral}~; le granite est une \omterm{def:g3:stones-pebbles-sand:rock}{roche}. Un \omterm{def:g6:rocks-and-minerals:mineral}{minéral} est une
seule substance naturelle dont les atomes sont disposés de façon
régulière~; une \omterm{def:g3:stones-pebbles-sand:rock}{roche} est faite d'un ou de plusieurs \omterm{def:g6:rocks-and-minerals:mineral}{minéraux} assemblés
(le granite~: feldspath, quartz et mica).
\end{solution}

\begin{solution}{exo:g6:rocks-and-minerals:3}
Talc (1), calcite (3), quartz (7), diamant (10). Le diamant, le plus dur,
raye tous les autres.
\end{solution}

\begin{solution}{exo:g6:rocks-and-minerals:4}
«~Une \omterm{def:g3:stones-pebbles-sand:rock}{roche} fondue devient une \omterm{def:g6:rocks-and-minerals:rock-families}{roche ignée} par refroidissement. Un calcaire
devient un marbre par la chaleur et la pression. Un granite se transforme
en grains de \omterm{def:g3:stones-pebbles-sand:sand}{sable} par \omterm{def:g5:water-shapes-the-land:weathering}{altération}.~»
\end{solution}

\begin{solution}{exo:g6:rocks-and-minerals:5}
(a) Basalte. (b) Marbre. (c) Grès.
\end{solution}

\begin{solution}{exo:g6:rocks-and-minerals:6}
Sa \omterm{def:g6:rocks-and-minerals:hardness}{dureté} est supérieure à celle d'un ongle (environ 2.5) et inférieure à
celle d'une pièce de cuivre (environ 3.5). Ce pourrait être de la calcite
(3).
\end{solution}

\begin{solution}{exo:g6:rocks-and-minerals:7}
Quartz~: $21 \div 60 = 0.35$, soit 35~\%. Feldspath~: $33 \div 60 = 0.55$,
soit 55~\%. Mica~: $6 \div 60 = 0.10$, soit 10~\%. Vérification~:
$35 + 55 + 10 = 100$.
\end{solution}

\begin{solution}{exo:g6:rocks-and-minerals:8}
Le granite refroidit très lentement, en profondeur~: ses grains ont
longtemps pour grandir, et ils deviennent gros. Le basalte vient de la
\omterm{def:g4:volcanoes-earthquakes-first-look:eruption}{lave}, qui refroidit vite en surface~: ses grains n'ont pas le temps de
grandir et restent trop petits pour être vus.
\end{solution}

\begin{solution}{exo:g6:rocks-and-minerals:9}
(a) Le calcaire est sédimentaire~; le marbre est métamorphique. (b) En
profondeur, la chaleur et la pression l'ont transformé sans le faire
fondre~: sa calcite a grandi en cristaux brillants qui s'emboîtent. (c) Tous
deux sont faits du même \omterm{def:g6:rocks-and-minerals:mineral}{minéral}, la calcite, et c'est la calcite qui fait
des bulles sous l'acide.
\end{solution}

\begin{solution}{exo:g6:rocks-and-minerals:10}
(a) Oui~: le quartz (7) est plus dur que le verre (environ 5.5). (b) Non~:
la fluorine (4) est plus tendre que le verre. (c) Oui~: un clou en acier
(environ 6.5) est plus dur que le feldspath (6).
\end{solution}

\begin{solution}{exo:g6:rocks-and-minerals:11}
Par exemple~: (1) le granite de la \omterm{def:g1:landscapes-around-us:landform}{montagne} (\omterm{def:g6:rocks-and-minerals:rock-families}{roche ignée}) est fragmenté par
l'\omterm{def:g5:water-shapes-the-land:weathering}{altération}, et ses grains de quartz sont libérés~; (2) une \omterm{def:g1:landscapes-around-us:water}{rivière} les
emporte~: c'est l'\omterm{def:g5:water-shapes-the-land:erosion}{érosion}~; (3) les grains se déposent en couches au fond
de la \omterm{def:g2:water-on-earth:ocean}{mer}~: c'est le \omterm{def:g5:water-shapes-the-land:deposition}{dépôt}~; (4) enfouis sous d'autres couches, ils sont
tassés et cimentés en un grès~: c'est le durcissement (la \omterm{def:g3:stones-pebbles-sand:rock}{roche} est
maintenant sédimentaire)~; (5) entraîné en profondeur, le grès est
transformé par la chaleur et la pression en une \omterm{def:g6:rocks-and-minerals:rock-families}{roche métamorphique}.
\end{solution}

\begin{solution}{exo:g6:rocks-and-minerals:12}
(a) Pas rayé par l'acier (environ 6.5) et capable de rayer le verre~: plus
dur que 6.5. Rayé par la topaze (8)~: plus tendre que 8. Sa \omterm{def:g6:rocks-and-minerals:hardness}{dureté} est 7~:
c'est du quartz. (b) Le quartz est plus dur que presque tout ce qui frotte
contre lui (autres grains, \omterm{def:g3:stones-pebbles-sand:pebble}{galets})~: il s'use à peine. La calcite (3) et le
mica (2 à 3) sont bien plus tendres~: le frottement les réduit en
minuscules morceaux qui sont emportés.
\end{solution}

\begin{solution}{pb:g6:rocks-and-minerals:1}
\textbf{1.} Du granite, une \omterm{def:g6:rocks-and-minerals:rock-families}{roche ignée}~; ses \omterm{def:g6:rocks-and-minerals:mineral}{minéraux} sont le feldspath
(rose ou blanc), le quartz (gris, vitreux) et le mica (noir, brillant).

\textbf{2.} Du calcaire, une \omterm{def:g6:rocks-and-minerals:rock-families}{roche sédimentaire}. Les coquilles
appartenaient à des animaux qui vivaient dans la \omterm{def:g2:water-on-earth:ocean}{mer}~; leurs restes se
sont déposés puis ont durci pour former la \omterm{def:g3:stones-pebbles-sand:rock}{roche}.

\textbf{3.} De l'ardoise, une \omterm{def:g6:rocks-and-minerals:rock-families}{roche métamorphique}, formée à partir d'argilite.

\textbf{4.} Du marbre, une \omterm{def:g6:rocks-and-minerals:rock-families}{roche métamorphique}, formé à partir de calcaire.

\textbf{5.} Du quartz~: pas rayé par l'acier (environ 6.5), il raye le
verre (environ 5.5)~; sa \omterm{def:g6:rocks-and-minerals:hardness}{dureté} est 7.

\textbf{6.} Rayés par l'acier (environ 6.5) mais pas par le verre (environ
5.5)~: entre 5.5 et 6.5. C'est du feldspath (orthose, 6).

\textbf{7.} Du mica (\omterm{def:g6:rocks-and-minerals:hardness}{dureté} 2 à 3).

\textbf{8.} Oui~: le pavage est en calcaire, fait de calcite (3), plus
tendre qu'une pièce de cuivre (environ 3.5).

\textbf{9.} L'\omterm{def:g5:water-shapes-the-land:weathering}{altération}~: la \omterm{def:g4:clouds-rain-snow:precipitation}{pluie} use lentement la surface de la pierre.

\textbf{10.} La \omterm{def:g1:landscapes-around-us:landform}{montagne} est altérée~; ses grains de quartz sont emportés
par les \omterm{def:g1:landscapes-around-us:water}{rivières} et se déposent en couches, au fond d'un \omterm{def:g1:landscapes-around-us:water}{lac} ou de la
\omterm{def:g2:water-on-earth:ocean}{mer}~; enfouis et durcis, ils peuvent devenir un grès, une \omterm{def:g6:rocks-and-minerals:rock-families}{roche
sédimentaire}.

\textbf{11.} Les trois familles sont présentes~: ignée (le granite de la
fontaine), sédimentaire (le calcaire du pavage), métamorphique (l'ardoise
du toit et le marbre de la statue).

\textbf{12.} Le \omterm{def:g6:rocks-and-minerals:mineral}{minéral} le plus dur est le quartz du granite~: \omterm{def:g6:rocks-and-minerals:hardness}{dureté} de
Mohs \textbf{7}.
\end{solution}
